\section{Introduction}

Vision-Language-Action (VLA) models such as $\pi_0$~\cite{black2024pi0}, OpenVLA~\cite{kim2024openvla}, and RT-2~\cite{brohan2023rt2} combine visual perception, language understanding, and action generation in a unified architecture to follow diverse natural language instructions across manipulation tasks. These models learn action policies from large-scale demonstration datasets and can generalize to new instructions, objects, and scenes. Yet when deployed in cluttered environments where non-target objects densely surround the manipulation target, VLAs frequently produce trajectories that collide with surrounding obstacles. The demonstration data used for training typically features unobstructed workspaces, and the resulting policies have no explicit incentive to maintain clearance from nearby surfaces. Collisions cause task failure, risk damage to the robot and surrounding objects, and represent a critical barrier to safe deployment in household and industrial settings. The SafeLIBERO benchmark~\cite{hu2025vlsa}, comprising 32 collision scenarios across 1,600 evaluation episodes, was introduced to systematically quantify this failure mode and confirms that even high-performing VLAs exhibit substantial collision rates in constrained workspaces. The root cause is architectural: current VLAs observe scenes through 2D images, generate action sequences through learned policies, and have no internal mechanism to determine whether the predicted end-effector trajectory will intersect obstacle surfaces.

Recent work addresses this problem, but exclusively through mechanisms external to the VLA's internal representations. AEGIS~\cite{hu2025vlsa} introduces a plug-in safety constraint layer based on Control Barrier Functions (CBFs): it identifies obstacles via open-set detection and depth estimation, then corrects unsafe actions after they have been generated, and achieves a Collision Avoidance Rate (CAR) of 77.85\% on SafeLIBERO. SafeVLA~\cite{zhang2025safevla} formulates safety as a constrained Markov Decision Process and reduces safety violations by 83.58\% through reward shaping during reinforcement learning; this approach, however, does not incorporate geometric information about the workspace and relies entirely on learned reward signals to encode safety. CoFreeVLA~\cite{cofreeevla2026} trains a risk estimator for dual-arm settings that detects dangerous configurations and halts execution when predicted collision risk exceeds a threshold. These approaches share a common architectural pattern: the safety mechanism operates independently of the VLA's learned feature representations. The model first generates a trajectory without collision awareness, and a separate module then evaluates or corrects it. This reactive paradigm faces two fundamental limitations. In cluttered environments, the space available for post-hoc trajectory correction shrinks, and corrections may push the end-effector into other obstacles. Because the VLA itself never learns to represent collision geometry, it cannot anticipate obstacles during action generation and will repeatedly propose unsafe trajectories that require external intervention. A proactive alternative would embed collision-relevant geometric information directly into the VLA's forward pass, so that action tokens encode collision awareness from the point of generation.

A parallel research direction has enriched VLA architectures with 3D spatial information. Over fourteen recent methods integrate depth maps, point clouds, or implicit 3D representations into VLA training or inference: DepthVLA~\cite{depthvla2025}, QDepth-VLA~\cite{li2025qdepth}, Spatial Forcing~\cite{li2025spatialforcing}, 3D-CAVLA~\cite{3dcavla2025}, GeoAware-VLA~\cite{geoaware2025}, SpatialVLA~\cite{spatialvla2025}, GST-VLA~\cite{gstvla2026}, StemVLA~\cite{stemvla2026}, ST-VLA~\cite{stvla2026}, Evo-0~\cite{evo02025}, SA-VLA~\cite{savla2026}, AugVLA-3D~\cite{augvla3d2026}, and Evo-Depth~\cite{evodepth2026}, among others. These methods improve task success rate, viewpoint generalization, or grasping precision, but none targets collision avoidance as a design objective and none evaluates on collision-specific metrics. QDepth-VLA, the closest methodological precedent for auxiliary spatial supervision in VLAs, uses quantized depth prediction as a training-time auxiliary task via a VQ-VAE decoder. Depth, however, is a 2.5D camera-frame representation: it encodes the distance from each pixel to the camera, and its gradients reside in image space. Deriving collision-relevant directions in the robot's workspace from these gradients requires an implicit perspective-to-workspace coordinate transformation that depends on camera intrinsics and extrinsics, and occluded regions (where obstacles may be hidden behind the robot arm or other objects) produce no depth signal. The result is a clear gap between two active communities: spatially-aware VLAs do not address collision avoidance, and collision-safe VLAs do not internalize spatial geometry into the model's representations.

We propose ColliForce-VLA, which integrates a differentiable workspace Signed Distance Field (SDF) into the VLA forward pass for collision-aware action generation. Workspace SDF encodes, at each 3D point, the signed distance to the nearest obstacle surface; its spatial gradient points directly away from that surface in workspace coordinates. Unlike depth gradients that reside in camera frame, SDF gradients require no coordinate transformation and cover the full 3D workspace including occluded regions. SDF gradients indicate repulsion direction from obstacles by construction, and their continuous nature provides informative signals even at moderate distances from surfaces, in contrast to binary occupancy whose gradient is zero almost everywhere. ColliForce-VLA augments $\pi_0$~\cite{black2024pi0} with two lightweight modules attached to the action expert: an SDF Decoder (4-layer MLP, ${\sim}$395K parameters) that predicts workspace SDF from the action expert's pooled suffix features, and an End-Effector (EE) Extractor (2-layer MLP, ${\sim}$263K parameters) that predicts end-effector position from the same features. During training, for near-collision samples where predicted SDF at the end-effector is low, we compute $\nabla_z \Phi$, the gradient of predicted SDF with respect to the shared feature representation $z$, and perform gradient ascent ($z' = z + \alpha \cdot \nabla_z \Phi$) to push $z$ toward collision-free configurations. Both $z$ and $z'$ are projected into action token space, and a knowledge distillation loss trains LoRA parameters (${\sim}$33.5M) to reproduce the refined actions. At inference, only the base model with trained LoRA is active; neither auxiliary module participates, and the method incurs zero additional latency. Unlike AEGIS~\cite{hu2025vlsa}, which corrects trajectories after generation, ColliForce-VLA shapes the action generation process through internalized collision geometry. Unlike QDepth-VLA~\cite{li2025qdepth}, which provides camera-frame auxiliary supervision only at training time without influencing action token generation, our SDF gradient refinement operates in workspace frame and modulates the feature representations from which action tokens are decoded. The collision-avoidance capability is distilled into the model's own LoRA parameters, not maintained by an external safety layer.

Our contributions are as follows:
\begin{itemize}
    \item We introduce a differentiable collision-aware mechanism for VLAs that integrates workspace SDF into the forward pass. SDF gradients in feature space ($\nabla_z \Phi$) guide action token generation toward collision-free configurations and provide a proactive alternative to post-hoc safety correction (Section~\ref{sec:method}).
    \item We design a knowledge distillation training strategy where an SDF gradient refiner serves as teacher and LoRA parameters as student. The model internalizes collision-avoidance behavior during training with zero additional inference cost (Section~\ref{sec:training}).
    \item We conduct systematic experiments and ablation studies on the SafeLIBERO benchmark that separately quantify the contributions of SDF auxiliary supervision and SDF gradient-guided refinement to Collision Avoidance Rate (CAR) and Task Success Rate (TSR) \textbf{[TODO: specific numbers pending experiments]} (Section~\ref{sec:experiments}).
\end{itemize}
